% -*- Mode:TeX -*-

%% The documentclass options along with the pagestyle can be used to generate
%% a technical report, a draft copy, or a regular thesis.  You may need to
%% re-specify the pagestyle after you \include  cover.tex.  For more
%% information, see the first few lines of mitthesis.cls. 

%\documentclass[12pt,vi,twoside]{mitthesis}
%%
%%  If you want your thesis copyright to you instead of MIT, use the
%%  ``vi'' option, as above.
%%
%\documentclass[12pt,twoside,leftblank]{mitthesis}
%%
%% If you want blank pages before new chapters to be labelled ``This
%% Page Intentionally Left Blank'', use the ``leftblank'' option, as
%% above. 

\documentclass[12pt,vi,twoside]{mitthesis}
\usepackage{lgrind}
\pagestyle{plain}
\usepackage[acronym,nonumberlist]{glossaries}
\usepackage{hyperref}


\usepackage{graphicx} % support the \includegraphics command and options
\usepackage[usenames,dvipsnames]{xcolor}
\usepackage{setspace}
\usepackage{amsmath}
\usepackage{mathtools}

\usepackage[parfill]{parskip} % Activate to begin paragraphs with an empty line rather than an indent
\setlength{\parindent}{1cm}
%%% PACKAGES
\usepackage{booktabs} % for much better looking tables
\usepackage{array} % for better arrays (eg matrices) in maths
\usepackage{paralist} % very flexible & customisable lists (eg. enumerate/itemize, etc.)
\usepackage{verbatim} % adds environment for commenting out blocks of text & for better verbatim
\usepackage{subfig} % make it possible to include more than one captioned figure/table in a single float
% These packages are all incorporated in the memoir class to one degree or another...

\usepackage{algorithmic}
\usepackage[chapter]{algorithm}
\usepackage{ifthen}



\usepackage{float}

\usepackage{tikz}
\usetikzlibrary{chains,fit,shapes}
\usetikzlibrary{arrows,positioning}
\usetikzlibrary{decorations.markings}



\usepackage{wrapfig}
\usepackage{listings}

 
\definecolor{dkgreen}{rgb}{0,0.6,0}
\definecolor{gray}{rgb}{0.5,0.5,0.5}
\definecolor{mauve}{rgb}{0.58,0,0.82}
 
\lstset{ %
  language=C,                % the language of the code
  basicstyle=\footnotesize,           % the size of the fonts that are used for the code
  %numbers=left,                   % where to put the line-numbers
  numberstyle=\tiny\color{gray},  % the style that is used for the line-numbers
  stepnumber=2,                   % the step between two line-numbers. If it's 1, each line 
                                  % will be numbered
  %numbersep=5pt,                  % how far the line-numbers are from the code
  backgroundcolor=\color{white},      % choose the background color. You must add \usepackage{color}
  showspaces=false,               % show spaces adding particular underscores
  showstringspaces=false,         % underline spaces within strings
  showtabs=false,                 % show tabs within strings adding particular underscores
  %frame=single,                   % adds a frame around the code
  rulecolor=\color{black},        % if not set, the frame-color may be changed on line-breaks within not-black text (e.g. commens (green here))
  tabsize=2,                      % sets default tabsize to 2 spaces
  captionpos=b,                   % sets the caption-position to bottom
  breaklines=true,                % sets automatic line breaking
  breakatwhitespace=false,        % sets if automatic breaks should only happen at whitespace
  title=\lstname,                   % show the filename of files included with \lstinputlisting;
                                  % also try caption instead of title
  keywordstyle=\color{blue},          % keyword style
  commentstyle=\color{dkgreen},       % comment style
  stringstyle=\color{mauve},         % string literal style
  escapeinside={\%*}{*)},            % if you want to add a comment within your code
  morekeywords={ushort},
  emph={thrust,device_ptr,cudaDeviceSynchronize,sort\_by\_key},emphstyle=\bf
  }


%%% HEADERS & FOOTERS
\usepackage{fancyhdr} % This should be set AFTER setting up the page geometry
\pagestyle{fancy} % options: empty , plain , fancy
\renewcommand{\headrulewidth}{0pt} % customise the layout...
\lhead{}\chead{}\rhead{}
\lfoot{}\cfoot{\thepage}\rfoot{}



%%% SECTION TITLE APPEARANCE
%\usepackage{sectsty}
%\allsectionsfont{\sffamily\mdseries\upshape} % (See the fntguide.pdf for font help)
% (This matches ConTeXt defaults)
%\renewcommand{\thesection}{\thesection}{\vspace{-1pt}}
%\usepackage[compact]{titlesec}
%\titlespacing{\section}{0pt}{*0}{*-1}
%\titlespacing{\subsection}{0pt}{*0}{*-1}
%\titlespacing{\subsubsection}{0pt}{*0}{*-1}


%%% ToC (table of contents) APPEARANCE
\usepackage[nottoc,notlof,notlot]{tocbibind} % Put the bibliography in the ToC
%\usepackage[titles,subfigure]{tocloft} % Alter the style of the Table of Contents
%\renewcommand{\cftsecfont}{\rmfamily\mdseries\upshape}
%\renewcommand{\cftsecpagefont}{\rmfamily\mdseries\upshape} % No bold!



%% This bit allows you to either specify only the files which you wish to
%% process, or `all' to process all files which you \include.
%% Krishna Sethuraman (1990).

%\typein [\files]{Enter file names to process, (chap1,chap2 ...), or `all' to
%process all files:}
%\def\all{all}
%\ifx\files\all \typeout{Including all files.} \else \typeout{Including only \files.} \includeonly{\files} \fi


\makeglossaries

%%%%%%%%%%%%%%%%%%%%%%%%%
% Acronyms
%%%%%%%%%%%%%%%%%%%%%%%%%

\newacronym{ac:ecc}{ECC}{Error-Correction Coding}

\newacronym{ac:mpi}{MPI}{Message Passing Interface}

\newacronym{ac:pic}{PIC}{Particle-In-Cell}

\newacronym{ac:gpu}{GPU}{Graphical Processing Unit}

\newacronym{ac:cpu}{CPU}{Central Processing Unit}

\newacronym{ac:cudpp}{CUDPP}{Cuda data parallel primitives}

\newacronym{ac:sm}{SM}{Streaming Multiprocessor}

\newacronym{ac:simd}{SIMD}{Single Instruction Multiple Data}
%%%%%%%%%%%%%%%%%%%%%%%%%
% Other Terms
%%%%%%%%%%%%%%%%%%%%%%%%%
\newglossaryentry{gls:CUDA}
{
  name=CUDA,
  description={NVIDIA's Compute Unified Architecture,
				 a general purpose computing architecture for NVIDIA GPUs that uses C as a high-level programming language}
}

\newglossaryentry{gls:atomicmemoryoperation}
{
  name=atomic memory operation,
  description={Memory access that to a device DRAM or shared memory that is guaranteed to be seen by all threads. Atomic operations lock down a memory address, delaying accesses by any other threads until the operation is complete.},
	plural=atomic memory operations
}

\newglossaryentry{gls:warp}
{
  name=warp,
  description={Grouping of threads on the GPU. One warp consists of 32 threads that execute simultaneously on a single \gls{gls:streamingmultiprocessor}. A warp executes one common instruction at a time, although not all threads in the warp need act on the issued instruction.},
	plural=warps
}

\newglossaryentry{gls:streamingmultiprocessor}
{
  name=streaming multiprocessor,
  description={SIMD processing unit on an NVIDIA GPU. Each streaming multiprocessor (SM) contains a number of streaming processors, also known as ``CUDA cores". The SM creates, manages, schedules and executes multiple \glspl{gls:warp} at a time. Many GPUs contain multiple SMs.\cite{NVIDIACorporation2011} }
}

\newglossaryentry{gls:threadblock}
{
  name=thread block,
  description={Grouping of threads on the GPU. Threads within the same block can communicate using the \_\_shared\_\_ memory space, and can be synchronized within a kernel using the \_\_syncthreads(); function. Thread blocks exist on a single SM, and can contain anywhere from 1 to 1024 threads, which are divided up into warps for execution. For additional information please see the \emph{NVIDA CUDA-C Programming Guide}\cite{NVIDIACorporation2011} }
}

\newglossaryentry{gls:kernel}
{
  name=kernel,
  description={CUDA function that is called by the CPU and executes on the GPU. When a kernel is launched multiple thread blocks are created and executed on the GPU. Kernels execute asynchronously with respect to the CPU code. Multiple kernels can exist simultaneously on the same GPU, and can be called asynchronously with memory transfers between CPU and GPU memory spaces. Kernels are also known as \_\_global\_\_ functions.}
}

\newglossaryentry{gls:host}
{
  name=host,
  description={A system consisting of the CPU and its memory space.}
}

\newglossaryentry{gls:device}
{
  name=device,
  description={A system consisting of a single GPU.}
}

\newglossaryentry{gls:sharedmemory}
{
  name=shared memory,
  description={User managed cache on a streaming multiprocessor. Shared memory can be nearly as fast as registers and can be used to facilitate communication between threads in the same thread block.}
}

\newglossaryentry{gls:globalmemory}
{
  name=global memory,
  description={GPU DRAM. Large memory space on the GPU, facilitates communication between threads in different thread blocks, and can be accessed by the CPU using cudaMemcpy() operations.}
}






\newglossaryentry{gls:thrust}
{
  name=THRUST,
  description={Parallel algorithm template library for GPUs and multi-core CPUs.\cite{NVIDIACorporation2011a} }
}


\newglossaryentry{gls:CUBLAS}
{
  name=CUBLAS,
  description={Fully GPU accelerated BLAS package provided with the CUDA toolkit.\cite{NVIDIACorporation2012} }
}

\newglossaryentry{gls:CUFFT}
{
  name=CUFFT,
  description={Fully GPU accelerated fast Fourier transform package provided with the CUDA toolkit.\cite{NVIDIACorporation2012a}}
}

\newglossaryentry{gls:CUSPARSE}
{
  name=CUSPARSE,
  description={Fully GPU accelerated sparse matrix linear algebra package provided with the CUDA toolkit.\cite{NVIDIACorporation2012c}}
}

\newglossaryentry{gls:CURAND}
{
  name=CURAND,
  description={Fully GPU accelerated, high quality pseudo-random and quasi-random number generation library provided with the CUDA toolkit.\cite{NVIDIACorporation2012b}}
}






























\renewcommand*{\glossaryname}{Frequently Used Terms}



\begin{document}

\include{cover}
\pagestyle{plain}
\include{contents}
\glsaddall
\printglossaries

%%%%%%%%%%%%%%%%%%%%%%%%%%%%%%%%%%%
\chapter{Introduction}
\label{ch:introduction}
%%%%%%%%%%%%%%%%%%%%%%%%%%%%%%%%%%%
	Over the past century humanity has become increasingly dependent on the 4th state of matter, plasma. Attaining a better understanding of plasma behaviour and interaction is critical to developing faster computer chips, creating new sources of energy, and expanding humanities influence among the stars. One important subset of plasma behaviour is how plasmas interact with solid objects such as dust particles, probes, and bodies traveling through space. These interactions can be very difficult to explore experimentally, and therefore must be modeled. 
	
A plasma's behaviour is heavily influenced by the collective electric and magnetic fields generated by the individual particles that comprise the plasma. This means that plasma behaviour is essentially a very large n-body problem, where for moderately dense plasmas n can be on the order of $10^{20}$ per cubic meter. No computer in existence can store the information for $10^{20}$ particles per cubic meter for large geometries, and calculating the interactions of every particle in the set with every other particle would be prohibitively long. The solution to this problem is to model only a subset of the particles. The modeled behaviour of these particles and their contributions to magnetic and electric fields can be used to statistically infer the behaviour of the rest of the plasma, essentially from first principles. This method is called particle-in-cell (PIC), and operates by moving particles on a potential grid and updating that potential with the new particle density at every timestep. 

The PIC method begins with an initial particle distribution, $\mathcal{P}(\mathbf{x}_i,\mathbf{v}_i)$. which is interpolated to the simulation mesh vertices in order to generate the charge and current density, $\mathcal{V}(\rho_j,\mathbf{J}_j)$. Integrating the charge and current densities defines the electric and magnetic fields, $\mathcal{V}(\mathbf{E}_j,\mathbf{B}_j)$ at the mesh vertices. From here the force, $\mathbf{F}_i$ on each particle is interpolated from $\mathbf{E}_j$ and $\mathbf{B}_j$. The force is then used to advance the positions and velocities of every particle. Following the advancing step, new particles must be reinjected to replace those that left the simulation domain. Collisions can be included after the reinjection step, although not all codes choose to include them. Finally new particle distribution function is interpolated to the mesh and the entire process starts again. Figure \ref{fig:pic_flowchart} shows the typical flow for this kind of electromagnetic PIC code.\cite{Verboncoeur2005}
\begin{figure}
\begin{center}
%\documentclass[a4paper,5pt]{article}
%\usepackage{tikz}
%\usetikzlibrary{arrows,positioning} 


%\begin{document}

\noindent
\begin{tikzpicture}[node distance=5cm,font=\footnotesize,auto]
\tikzset{
    %Define standard arrow tip
    >=stealth',
    %Define style for boxes
    punkt/.style={
           rectangle,
           rounded corners,
           draw=black, very thick,
           text width=7.0em,
			  text height=0.5em,
			  minimum width=7.0em,
           minimum height=4.0em,
			  },
    % Define arrow style
    pil/.style={
           ->,
           very thick,
           shorten <=2pt,
           shorten >=2pt,}
}
	\noindent
 	\node[shift={(5cm,0cm)}] (dtcenter) at (0,0) {$\Delta t$};
	\node[shift={(160:1cm)}] (circstart) at (dtcenter) {};
	\draw[pil] (circstart) arc (160:-160:1cm);
 %nodes

 	\node[punkt,shift={(-2.5cm,2.5cm)}] (padvnc) at (dtcenter) {Integrate equations of motion\newline $\mathbf{F}_i \rightarrow \mathbf{v}_i' \rightarrow \mathbf{x}_i$};

 	\node[punkt,shift={(2.5cm,2.5cm)}] (ploss) at (dtcenter) {Particle loss and re-injection};

	\node[punkt,shift={(7.0cm,0)}] (collide) at (dtcenter) {Monte Carlo collisions\newline $\mathbf{v}_i' \rightarrow \mathbf{v}_i$};

	\node[punkt,shift={(2.5cm,-2.5)}] (c2mesh) at (dtcenter) {Interpolate particles to mesh\newline $(\mathbf{x}_i,\mathbf{v}_i) \rightarrow \rho_j$};

	\node[punkt,shift={(-2.5cm,-2.5)}] (psolve) at (dtcenter) {Integrate field equations\newline $\rho_j \rightarrow \mathbf{E}_j$};

	\node[punkt,shift={(-5.5cm,0)}] (mesh2p) at (dtcenter) {Interpolate field to particles\newline $\mathbf{E}_j \rightarrow \mathbf{F}_i$};

	\path (padvnc.east) edge[pil] (ploss.west);
	\path (ploss.east) edge[very thick,dashed,out=-60, in=60,shorten <=2pt,shorten >=2pt,] (c2mesh.east);
	\path (ploss.east) edge[pil,out=0, in=90] (collide.north);
	\path (collide.south) edge[pil,out=-90, in=0] (c2mesh.east);
	\path (c2mesh.west) edge[pil] (psolve.east);
	\path (psolve.west) edge[pil,out=180,in=-90] (mesh2p.south);
	\path (mesh2p.north) edge[pil,out=90,in=180] (padvnc.west);
	

\end{tikzpicture}


%\end{document}

%%% Local Variables:
%%% mode: latex
%%% TeX-master: t
%%% End:

\end{center}
\caption[Flow schematic for the PIC method.]{Flow schematic for the electromagnetic PIC method. Monte Carlo collisions are optional and not included in all codes.}
\label{fig:pic_flowchart}
\end{figure}


%%%%%%%%%%%%%%%%%%%%%%%%%%%%%%%%%%%% 
	\section{Motivation}
%%%%%%%%%%%%%%%%%%%%%%%%%%%%%%%%%%%%

	The PIC method is very good at modeling complicated plasma behaviour, however this method relies on tracking very large numbers of particles for good statistics. In order to achieve ``good'' statistics PIC codes employ millions to billions of particles, which means that these codes require a very large amount of computation time for each timestep. Running billions of particles on a single core for thousands of timesteps is not really feasible, it simply takes too long to compute a solution. 
	
	One way to reduce the total run time of PIC codes is to parallelize them. Since PIC codes operate on the fact that the potential changes little over the course of a single timestep, each particle can be assumed to be independent of its neighbors. This leads to a situation that is trivially parallel. In theory a machine with a million processors could run every particle on a separate processor. This of course assumes that the majority of the computational complexity lies in moving the particles and that communication between processors is very fast.  


%%%%%%%%%%%%%%%%%%%%%%%%%%%%%%%%%%%%
		\subsection{GPUs vs CPUs}
%%%%%%%%%%%%%%%%%%%%%%%%%%%%%%%%%%%%
	The ideal computing system for a particle in cell code should have a large number of relatively simple processors with very low communication overhead. Traditional CPUs are just the opposite of this. CPUs tend to have 4-8 complicated processors that are very good at performing large operations on small sets of data, but very slow at communicating between multiple processors. CPUs are designed to be able to actively switch tasks on the fly, which makes them very good at simultaneously running web-browser, decoding a video, and playing a video game. However, this flexibility requires a large number of cycles to switch between tasks, and a large amount of cache to store partially completed tasks.

Graphical processing units, or GPUs, forgo the flexibility of CPUs in favor of raw processing capability. Reducing the size of the cache and employing single instruction multiple data (SIMD) parallelism allows GPU manufacturers to combine hundreds of processors on a single chip. In order to supply enough data to keep hundreds of processors busy, GPUs also have a very large data channel between processors and DRAM. All of these features are chosen to create a math processor that excels at operations where each processor operates on data that is invisible to the other processors. These features give GPUs a significant raw floating point performance advantage over CPUs, as shown in figure \ref{fig:gpu_vs_cpu}. 


\begin{figure}
\begin{center}
\includegraphics[width=5in]{introduction/gpu_vs_cpu.png}
\end{center}
\caption[Performance comparison of GPUs vs CPUs.]{Performance comparison of GPUs vs CPUs. GPU performance is continuing to increase at a very rapid pace.\cite{NVIDIACorporation2011}}
\label{fig:gpu_vs_cpu}
\end{figure}

The hardware in GPUs is tailored to excel at performing tasks such as ray-tracing, which is very similar to particle moving. Therefore it is by no means unreasonable to conclude that GPUs can be very good PIC code processors. The advantages that GPUs have over CPUs for scientific computing include:

\begin{itemize}
	\item Higher performance per unit cost.
	\item Higher performance per watt.
	\item Easier to upgrade.
	\item GPUs are still improving with Moore's law.
\end{itemize}

All of their advantages are observed when comparing the CPU and GPU versions of the same PIC code. While these advantages are very promising there are also several disadvantages to GPU computing:

\begin{itemize}
	\item Increased code complexity.
	\item Smaller memory space.
	\item Smaller cache.
	\item Slow communication between CPU and GPU.
	\item Algorithms dependent on hardware configuration.
\end{itemize}

The key to developing efficient PIC algorithms that utilize GPUs lies in balancing the work between the two architectures. Some operations will be easier to implement on the CPU and be just as fast as the GPU, while others will be significantly faster on the GPU. Partitioning the code between different architectures outlines a very important aspect of parallel computing; multiple levels of parallelism.

%%%%%%%%%%%%%%%%%%%%%%%%%%%%%%%%%%%%
	\section{Multiple Levels of Parallelism}
%%%%%%%%%%%%%%%%%%%%%%%%%%%%%%%%%%%%
	Currently most parallelization is done by dividing a task between a number of threads on different CPUs, and using an communication and synchronization API such as \gls{ac:mpi} to allow those threads to communicate. This network of threads has a master node, usually node 0, which orchestrates the communication between other nodes. This is analogous to how a single CPU-GPU system operates. The CPU is the ``Master'' and serves as a communication hub for groups of execution threads on the GPU called thread blocks. Each thread block is itself a cluster of threads that can communicate through a memory space aptly named ``shared memory''. When all of these systems are used together the resulting architecture has multiple levels of parallelism. The levels parallelism of a multi-GPU architecture are shown in figure \ref{fig:multiparallel}.

\begin{figure}
\begin{center}
\includegraphics[width=5in]{introduction/multi_parallel.png}
\end{center}
\caption[Multiple levels of parallelism.]{Multiple levels of parallelism. (1) Cluster of systems communicating through a LAN. (2) Multiple GPUs per system communicating through system DRAM. (3) Multiple streaming multiprocessors per GPU execute thread-blocks and communicate through GPU global memory. (4) Multiple \gls{gls:CUDA} cores per multiprocessor execute thread-warps and communicate through on chip shared memory. }
\label{fig:multiparallel}
\end{figure}

The point is that multiple domain decompositions must be performed in order to fully utilize the capabilities of this system. The coarse decomposition is very similar to that used for \gls{ac:mpi} systems, but the fine decomposition can be very different due to the significantly higher memory bandwidth and smaller cache of GPUs. 

%%%%%%%%%%%%%%%%%%%%%%%%%%%%%%%%%%%%
		\subsection{Parallelization Opportunities in PIC Codes}
%%%%%%%%%%%%%%%%%%%%%%%%%%%%%%%%%%%%

\begin{figure}
\begin{center}
%\documentclass[a4paper,10pt]{article}
%\usepackage[usenames,dvipsnames]{xcolor}
%\usepackage{tikz}
%\usetikzlibrary{arrows,positioning} 
%\usetikzlibrary{chains,fit,shapes}



%\begin{document}


\noindent
\begin{tikzpicture}[node distance=5cm,font=\footnotesize,auto]

\tikzset{
    %Define standard arrow tip
    >=stealth',
    %Define style for boxes
    punkt/.style={
           rectangle,
           rounded corners,
           draw=none, very thick,
           text width=2.5cm,
			  minimum width=2.6cm,
           minimum height=1.5cm,
			  },
    % Define arrow style
    pil/.style={
           ->,
           very thick,
           shorten <=2pt,
           shorten >=2pt,},
    punkt2/.style={
           rectangle,
           rounded corners,
           draw=black, very thick,
			  minimum width=2.65cm,
           minimum height=1.35cm,
			  anchor=west},
    punkt3/.style={
           rectangle,
           rounded corners=15pt,
           draw=black,
			  minimum width=3.2cm,
           minimum height=2.0cm,
			  }
}

	\noindent
 	\node[shift={(0cm,0cm)}] (dtcenter) at (0,0) {$\Delta t$};
	\node[shift={(160:1cm)}] (circstart) at (dtcenter) {};
	\draw[pil] (circstart) arc (160:-160:1cm);
 %nodes

 	\node[punkt,shift={(-2.5cm,2.5cm)}] (padvnc2) at (dtcenter) {Integrate equations of motion\newline $\mathbf{F}_i \rightarrow \mathbf{v}_i' \rightarrow \mathbf{x}_i$};

 	\node[punkt,shift={(2.5cm,2.5cm)}] (ploss2) at (dtcenter) {Particle loss and re-injection};

	\node[punkt,shift={(7.0cm,0)}] (collide2) at (dtcenter) {Monte Carlo collisions\newline $\mathbf{v}_i' \rightarrow \mathbf{v}_i$};

	\node[punkt,shift={(2.5cm,-2.5)}] (c2mesh2) at (dtcenter) {Interpolate particles to mesh\newline $(\mathbf{x}_i,\mathbf{v}_i) \rightarrow \rho_j$};

	\node[punkt,shift={(-2.5cm,-2.5)}] (psolve2) at (dtcenter) {Integrate field equations\newline $\rho_j \rightarrow \mathbf{E}_j$};

	\node[punkt,shift={(-5.5cm,0)}] (mesh2p2) at (dtcenter) {Interpolate field to particles\newline $\mathbf{E}_j \rightarrow \mathbf{F}_i$};

	\node[punkt2] (padvnc) at (padvnc2.west) {};
	\node[punkt2] (ploss) at (ploss2.west) {};
	\node[punkt2] (collide) at (collide2.west) {};
	\node[punkt2] (c2mesh) at (c2mesh2.west) {};
	\node[punkt2] (psolve) at (psolve2.west) {};
	\node[punkt2] (mesh2p) at (mesh2p2.west) {};

	\path (padvnc.east) edge[pil] (ploss.west);
	\path (ploss.east) edge[very thick,dashed,out=-60, in=60,shorten <=2pt,shorten >=2pt,] node[auto] {(2)} (c2mesh.east);
	\path (ploss.east) edge[pil,out=0, in=90] node[below] {(1)} (collide.north);
	\path (collide.south) edge[pil,out=-90, in=0] (c2mesh.east);
	\path (c2mesh.west) edge[pil] node[above] {Reduce $\rho_j$} (psolve.east) ;
	\path (psolve.west) edge[pil,out=180,in=-90] node[auto] {Bcast $\mathbf{E}_j$} (mesh2p.south);
	\path (mesh2p.north) edge[pil,out=90,in=180] (padvnc.west);

	
\begin{scope}[color=Emerald,very thick]
\newcommand*{\pshiftmm}{(-2.0mm,-2.0mm)}
\newcommand*{\pshiftpm}{(2.0mm,-2.0mm)}
\newcommand*{\pshiftmp}{(-2.0mm,2.0mm)}
\newcommand*{\pshiftpp}{(2.0mm,2.0mm)}

	\coordinate[shift=\pshiftpm] (p1) at (mesh2p.south east) {};
	\coordinate[shift=\pshiftpm] (p2) at (padvnc.south east) {};
	\coordinate[shift=\pshiftmm] (p3) at (ploss.south west) {};

	\coordinate[shift=\pshiftmm] (p4) at (collide.south west) {};
	\coordinate[shift=\pshiftpm] (p5) at (collide.south east) {};
	\coordinate[shift=\pshiftpp] (p6) at (collide.north east) {};

	\coordinate[shift=\pshiftpp] (p7) at (ploss.north east) {};

	\coordinate[shift=\pshiftmp] (p8) at (padvnc.north west) {};
	\coordinate[shift=\pshiftmp] (p9) at (mesh2p.north west) {};
	\coordinate[shift=\pshiftmm] (p10) at (mesh2p.south west) {};
	\coordinate[shift=\pshiftmm] (p11) at (mesh2p.south west) {};

	\coordinate[shift=\pshiftpp] (c1) at (mesh2p.north east) {};
	\coordinate[shift=\pshiftmm] (c2) at (padvnc.south west) {};
	\coordinate[shift=\pshiftpp] (c3) at (padvnc.south east) {};
	\coordinate[shift=\pshiftpp] (c4) at (ploss.south west) {};
	\coordinate[shift=\pshiftpm] (c5) at (ploss.south east) {};
	\coordinate[shift=\pshiftmp] (c6) at (collide.north west) {};

	\draw[rounded corners=10pt] (p11) -- (p1) .. controls (c1) and (c2) .. (p2) 
											-- node[auto] {(A)} (p3) .. controls (c5) and (c6) .. (p4) 
											-- (p5) 
											--  (p6) [bend left=30] -- (p7) --
											 (p8) [bend right=30] -- (p9) --
											(p10) -- cycle;
	\node[punkt3,color=Purple] (PP1) at (c2mesh.center) {};
	\node[above,color=Purple] at (PP1.north) {(B)};
	\node[punkt3,color=Maroon] (PP2) at (psolve.center) {};
	\node[above,color=Maroon] at (PP2.north) {(C)};
\end{scope}


	

\end{tikzpicture}


%\end{document}

%%% Local Variables:
%%% mode: latex
%%% TeX-master: t
%%% End:

\end{center}
\caption[Parallization opportunities in the PIC method]{Flow schematic for the electrostatic PIC method with parallelizable steps highlighted. (1) and (2) are possible paths that the code can take depending on whether or not collisions are used. (A) Particle steps are embarassingly parallel. (B) Particle to Mesh interpolation steps are difficult to parallelize. (C) Field solving step can make use of parallelized linear algebra libraries.}
\label{fig:pic_flowchart_parallel}
\end{figure}

Most of the steps in the PIC method can be parallelized, although some steps are more difficult than others. Figure \ref{fig:pic_flowchart_parallel} highlights the steps of the PIC method that are parallelizeable, and groups them according to how they can be parallelized. The particle steps, group (A) in the figure, are probably the easiest steps to implement in parallel. In theory every particle could be advanced by a separate processor or thread. In the case of GPUs it is entirely feasible to assign a single thread to every particle. There are several issues that arise when implementing steps that contain a large amount of execution divergence on the GPU. 
Execution divergence occurs whenever the computation path for two particles requires different instructions, such as if one particle must evaluate the contents of an \emph{if} statement, while another does not.
The particle loss / re-injection step, and the collision step are two such subroutines in which the execution paths of particles that are being lost, re-injected, or colliding will diverge greatly from the majority of the particles. In these cases it is best to operate on subsets of the particle list in which all particles of the subset follow the same execution path. There are ways in which subsets of the particle list can be operated on efficiently without large amounts of execution divergence. 

Parallelizing the particle to mesh step $((\mathbf{x}_i,\mathbf{v}_i)\rightarrow \rho_j)$, (B) in figure \ref{fig:pic_flowchart_parallel}, is trickier, depending on the amount of memory available per thread. In the case of multi-CPU implementations, each thread has enough memory to store a copy of the entire grid. This thread can process a particle that exists anywhere on the grid, and map its contribution to the grid without memory conflicts. The case is very different for GPUs, which have insufficient memory to support replicating the entire grid for every thread. Therefore the key to parallelizing the particle-to-mesh interpolation step on the GPU is determining way to efficiently avoid memory conflicts, using some kind of \gls{gls:atomicmemoryoperation} or domain decomposition. 

The last step (C), the field solve $(\rho_j \rightarrow \mathbf{E}_j)$, can be fairly straightforward to parallelize. The field solve involves integrating Maxwell's equations over the entire grid. In the case of electrostatic codes this is simply solving Poisson's equation. Poisson's equation is a straightforward linear equation that can be solved using a number of parallel $Ax=b$ solvers, popular ones being preconditioned bi-conjugate gradient solvers and spectral solvers.

The vast majority of the PIC algorithm is well suited for implementation on the GPU. There are some challenges to overcome, such as how to perform the charge assign, and how to efficiently handle re-injections and collisions. Currently GPU computing is in its infancy, and the solutions to these challenges are evolving rapidly. In the following section we will look at some of the solutions that others have developed to overcome these issues. 

%%%%%%%%%%%%%%%%%%%%%%%%%%%%%%%%%%%%
		\section{GPU PIC Code Development}
%%%%%%%%%%%%%%%%%%%%%%%%%%%%%%%%%%%%
	Before we dive into past research into implementing PIC on GPUs it would be useful to know a little more about general purpose graphical processor (GPGPU) computing. Currently there are two frameworks for writing GPGPU programs, the Open Computing Language (OpenCL) initially developed by Apple Inc., and the NVIDIA Corporation's Compute Unified Architecture \gls{gls:CUDA}.\cite{NVIDIACorporation2012d,NVIDIACorporation2011} Each framework has its own strengths and weaknesses. \gls{gls:CUDA} is a software environment that allows developers to use C as a high-level programming language. Developing in \gls{gls:CUDA} C is straightforward and is very well documented and supported by NVIDIA. The \gls{gls:CUDA} toolkit and SDK provided by NVIDIA include decent debugging and profiling tools as well as a large number of fully GPU accelerated libraries. These libraries include:

\begin{itemize}
\item \gls{gls:thrust} (template and algorithm library)\cite{NVIDIACorporation2011a}
\item Cula (linear algebra library)
\item CUBLAS (Basic Linear Algebra Subprograms)\cite{NVIDIACorporation2012}
\item CUFFT (Fast Fourier Transform (FFT) Library)\cite{NVIDIACorporation2012a}
\item CUSPARSE (Sparse linear algebra routines)\cite{NVIDIACorporation2012c}
\item CURAND (high-quality pseudo-random and quasi-random numbers)\cite{NVIDIACorporation2012b}
\end{itemize}

The primary downside to \gls{gls:CUDA} is that it is restricted to NVIDIA GPUs, and cannot be used for other hardware platforms. OpenCL on the other hand supports multiple computing platforms and has been adopted by Intel, Advanced Micro Devices, Nvidia, and ARM Holdings. However, developing in OpenCL can be somewhat cumbersome compared to \gls{gls:CUDA}, especially when it comes to debugging and profiling. OpenCL also lacks mature libraries similar to those provided with \gls{gls:CUDA}. For these reasons, and the fact that the primary compute-specific cards are the NVIDA Tesla cards, most GPU PIC code development has been done using \gls{gls:CUDA}.\cite{NVIDIACorporation2012d,NVIDIACorporation2011}

%%%%%%%%%%%%%%%%%%%%%%%%%%%%%%%%%%%%
		\section{Current Status of GPU PIC codes}
%%%%%%%%%%%%%%%%%%%%%%%%%%%%%%%%%%%%

						Some work on efficient GPU based PIC codes has already been done. This past work will be briefly introduced here and discussed in depth in chapter \ref{ch:design}. One of the earliest efforts in developing plasma PIC simulations on the GPU is that of George Stantchev et al. In 2008 Stantchev et al published a paper in the \emph{Journal of Parallel Distributed Computing}\cite{Stantchev2008} outlining the development of a fast density accumulation method for PIC implementations on the GPU. Stantchev's paper was one of the first to identify the issue of parallelizing the density update and develop a very efficient solution. The solution involves defining clusters of grid cells and sorting the particle list according to which cluster they belong to. The cell clusters are small enough to fit into \emph{\_\_shared\_\_} memory and updated pseudo atomically using a now obsolete thread tagging technique. The simplicity and efficiency of this method has led to its use in many other GPU PIC implementations. Stantchev also identified the issue of the particle sort step required by their density accumulation method. Developing an optimized sort has since become one of the biggest challenges for GPU PIC implementations. 

Also in 2008, Dominique Aubert et al published their development of a GPGPU galactic dynamics PIC code.\cite{Aubert2009}. Aubert et al adapted their entire code for GPUs and tested several field-solving and density-accumulation techniques. The first density accumulation technique relied on the radix sort provided by the \gls{gls:CUDA} Data Parallel Primitives (CUDPP) library to order the particle data and construct a histogram. The second technique used the GPU to find the nearest cells for each particle and then construct the density histogram on the CPU. The details concerning how these histograms were constructed are not disclosed in the paper. As for field solving techniques, GPU versions of both FFT and multi-grid Poisson solvers were used. The FFT solver made use of the CUFFT API provided by the \gls{gls:CUDA} Toolkit. The multi-grid (MG) solver was a from scratch GPU implementation of the MG solver outlined in \cite{NumericalRecipes}. Aubert et al achieved speed ups of 40x for the FFT and 60x for MG, and approximately 20x for the velocity and position updates, but the histogramming was on the order of 5-10x \emph{slower} on the GPU. They acknowledge that their histogram method is too slow, and propose using an improved technique following that developed by Stantchev et al.

A further step in GPU PIC codes was taken in 2010 by the simulation code PIConGPU\cite{Burau2010}, developed by Heiko Burau et al. PIConGPU was one of the first multi-GPU PIC implementations involving multiple nodes. PIConGPU is a fully relativistic, 2D electromagnetic code used to simulate the acceleration of electrons in an under-dense plasma by a laser-driven wakefield. Burau et al also identified the issue of parallelizing charge and current density accumulation on SIMD architectures and came up with a unique solution, a linked particle list. This linked particle list is one of the core features of PIConGPU and is geared towards maintaining localized data access for each kernel. In order to facilitate scalability over multiple nodes the simulation volume is domain decomposed and distributed across multiple GPUs. Each sub-domain is bordered by a region of guard cells to facilitate particle transfers between sub-domains. These guard cells overlap with adjacent domains and serve as a buffer for particles moving from one sub-domain to the next. The results of this research indicated that single GPU PIC algorithms can be scaled to multi-node clusters despite high host-device communication latency.

In March, 2010 NVIDIA released the first GPUs utilizing the Fermi architecture. Fermi boasted several key improvements for scientific computing over the G80 and GT200 architectures. These improvements included \cite{NVIDIACorporation2009}:

\begin{itemize}
\singlespacing
\item Configurable L1 and unified L2 caches.
\item 8x the double precision performance of the GT200 architecture. 
\item Error-Correction Code (ECC) support. 
\item Improved \gls{gls:atomicmemoryoperation} performance.
\item Full IEEE 754-2008 32-bit and 64-bit arithmetic floating point precision compliance. 
\end{itemize}

The significant increase in double precision performance and addition of ECC support are desirable for scientific computing on GPUs. With the launch of Fermi, work on GPU based PIC codes increased significantly. Another 2D relativistic GPU PIC code was developed by Paulo Abreu et al in late 2010.\cite{Abreu2011} Abreu et al chose an approach to the charge and current accumulation which was in some regards opposite of that employed by Stantchev et al. This approach was based on using pseudo \glspl{gls:atomicmemoryoperation} that actually consisted of two \textbf{atomicExchange}$()$ operations. In order to minimize atomic collisions the particle list is ordered such that the probability that two threads within a \gls{gls:warp} were operating on particles in the same cell is low. The initial particle distribution is prepared in way to minimize this probability, but this distribution can be upset throughout the simulation. If the distribution degradation reaches a certain threshold then a redistribution of the particle list takes place. This redistribution was handled by a sorting operation, using the CUDPP radix sort, and a \emph{stride} distance, or the distance in particle index between particles handled by one thread and the next. The resulting code had a relatively fast particle advance but the current deposition was rather slow. The current deposition step constituted approximately 67\% of the code runtime for 1.2 million particles on a $128^2$ grid.

One of the more unique GPU PIC implementations was developed by Rejith Joseph et al and published at the 2011 \emph{IEEE International Parallel \& Distributed Processing Symposium}.\cite{Joseph2011} This code, a GPU implementation of the XGC1 code differed from previous codes in two ways. First, XGC1 uses an asymmetric triangular mesh, unlike most GPU PIC codes which use symmetric rectangular meshes. Second, there are no limits imposed on particle movement, which makes the lessons taught by the code more applicable in general.

The first point, the use of a triangular mesh is interesting because most of the time a symmetric rectangular mesh is preferred on the GPU. Triangular meshes can be more efficient at simulating irregular domains by using finer spatial discretization for important regions.  As in other PIC codes, in order to facilitate a fast particle to grid interpolation it was important that all of the triangles that a thread-block might access fit into the limited \emph{\_\_shared\_\_} memory. This limit, imposed by the size of \emph{\_\_shared\_\_} memory, necessitated that the primary mesh be partitioned into smaller regions. Joseph et al tried three different partitioning strategies, density based, uniform, and non-uniform. The density based strategy partitioned the mesh such that each region had approximately the same triangle density. The uniform approach divided the mesh into large uniform rectangular regions. Lastly the non-uniform partitioning was a sort of hybrid of the uniform and density approaches, using rectangular regions that varied in size depending on the granularity of the triangle density.\cite{Joseph2011}

The goal of these partitioning schemes was to take care of load balance on the GPU. As it turns out, the automatic load balancing provided by the GPU hardware is very good. When Joseph et al compared triangle search times between the different partitioning schemes the uniform scheme was the fastest for all numbers of thread blocks, but for very large numbers of blocks the Uniform and the Non-Uniform schemes converged to roughly the same execution time. This brings up a very interesting point that can teach us a valuable lesson; namely, the GPU tends to like very simple problems and often times increasing the complexity results in reduced performance. While complicated performance enhancing techniques work well on serial architectures, often they will only hurt you when implemented on the GPU. 

The final code that we will look at was published in the \emph{Journal of Computational Physics} by Xianglong Kong et al in 2011.\cite{Kong2011} This code was a 2D3V fully relativistic electromagnetic code that achieved speed-up factors of 81x and 27x over the CPU runtime for cold plasma runs and extremely relativistic plasma runs respectively. This code approached the issue of maintaining an organized particle list very differently than previous codes. The approach used in this code is more akin to message passing than sorting, and will be explained in more detail in chapter \ref{ch:design}. The performance of this code greatly depended on the fraction of particles that crossed from one sub-region into the next, or the crossing fraction $\eta$. The achieved run-time per particle-step was 4.6 ns for $\eta=0.5\%$, and up to 9.15 ns for $\eta=7.5\%$. 

		

%%%%%%%%%%%%%%%%%%%%%%%%%%%%%%%%%%%%
	\section{SCEPTIC3D}
%%%%%%%%%%%%%%%%%%%%%%%%%%%%%%%%%%%%
Now that previous works in the realm of GPU PIC codes have been discussed, it is time to introduce SCEPTIC3D, the code that we chose to implement a GPU version of. SCEPTIC3D is three dimensional hybrid PIC code specifically designed to solve the problem of ion flow past a negatively biased sphere in a uniform magnetic field. The current version of the code was derived from the 2D/3v code SCEPTIC which was originally written by Ian Hutchinson \cite{Hutchinson2002,Hutchinson2003,Hutchinson2005,Hutchinson2006}. The code is written in FORTRAN, and has been previously parallelized using standard \gls{ac:mpi} communication interface. Uses of SCEPTIC3D include exploring probe-plasma interactions, modeling the behaviour of particle grains in dusty plasmas, and determining the interactions between objects in space and their environment. As previously mentioned, the PIC method requires a large number of particles in order to reduce statistical noise. In the case of SCEPTIC3D, runs typically consist of 50 million particles.
Moving 50 million particles on a $64^3$ grid takes roughly 11 seconds per step using 4 cores of an Intel i7 930 processor. The fact that SCEPTIC3D has already been parallelized using \gls{ac:mpi} makes it far easier to develop a multi-GPU implementation. 

	\subsection{CPU Code Profiling}


\begin{figure}
\begin{center}
\input{introduction/sceptic3D_profile.tex}
\end{center}
\caption[Breakdown of SCEPTIC3D runtime costs]{Breakdown of SCEPTIC3D runtime costs by subroutine. This is for 12.5 million particles on a $64^3$ grid. Times are in ns per particle step.}
\label{fig:sceptic3D_profile}
\end{figure}

In figure \ref{fig:sceptic3D_profile} we break down the runtime of SCEPTIC3D into the main subroutines. This figure indicates that porting the charge assign and particle advancing steps should be the main priority. Of course, if we hope to get a performance boost of more than 3x for both the particle advancing and charge assign steps, then we will have to worry about the field solve as well. However, it should be noted that the balance between the particle and field costs depends on the number of particles per cell. In the case of figure \ref{fig:sceptic3D_profile} we used roughly 48 particles per cell, generally a larger fraction is preferable in order to reduce particle noise. 

%\begin{figure}
%\begin{center}
%\includegraphics[width=4in]{introduction/not_finished.pdf}
%\end{center}
%\caption{Flow schematic for the PIC method with sceptic subroutine names Need to make figure}
%\label{fig:pic_flowchart_sceptic}
%\end{figure}


The GPU implementation of SCEPTIC3D is very different from any other GPU PIC implementation. While we will focus a great deal on performance, we will also place a great deal of emphasis on implementation difficulty and applicability to future PIC implementations on the GPU. We will determine how critical the performance of the field solver is when the run time of other steps is significantly reduced. Lastly, we will look at how primarily graphics related GPU features can be utilized by PIC codes. 














%% This is an example first chapter.  You should put chapter/appendix that you
%% write into a separate file, and add a line \include{yourfilename} to
%% main.tex, where `yourfilename.tex' is the name of the chapter/appendix file.
%% You can process specific files by typing their names in at the 
%% \files=
%% prompt when you run the file main.tex through LaTeX.
\chapter{Design Options}
\label{ch:design}

	\section{Particle List Structure}
		\subsection{Other Codes}
		\subsection{In house tests}
\noindent \begin{table}[h]
\begin{tabular}{| p{4.0cm} | p{3.5cm} | p{2.5cm} | p{4.0cm} |}
\hline
Component & SoA (ms) & AoS (ms) & Speedup (SoA vs AoS) \\ \hline
Particle data read, move, and write & 758 & 955 & 1.26x \\ \hline
Count Particles & 32.7 & 109 & 3.35x \\ \hline
Data Reorder & 346 & 480 & 1.38x \\ \hline
Total CPU run time & 2491 & 3284 & 1.31x \\ \hline
\end{tabular}
\caption{Execution times of main steps for Array of Structures and Structure of Arrays. Count Particles and Data Reorder are steps used for a sorted particle list. Count Particles counts the number of particles in each sub-domain. Data Reorder reorders the particle list data after the binindex / particle ID pair have been sorted by the radix sort.}
\label{tab:struct_compare} 
\end{table}

	\section{Charge Assign}
		\subsection{Naive Atomic Approach}
		\subsection{Other Codes}

	\section{Particle List Sort}
		\subsection{Costs and Benefits}
		\subsection{Other Codes}
		*Stantchev Particle Binning
		*Kong Particle Passing
		*Linked Particle List
		\subsection{In house tests}

	\section{Particle Advancing}
		\subsection{Assumptions}
		\subsection{Other Codes}
		\subsection{Reinjections and Diagnostics}

	\section{Poisson Solve}
		\subsection{Desired Performance}
		\subsection{Performance vs Implementation Difficulty}

	\section{Grid Dimension Constraints and Handling}

%% This is an example first chapter.  You should put chapter/appendix that you
%% write into a separate file, and add a line \include{yourfilename} to
%% main.tex, where `yourfilename.tex' is the name of the chapter/appendix file.
%% You can process specific files by typing their names in at the 
%% \files=
%% prompt when you run the file main.tex through LaTeX.
\chapter{Implementation} 
\label{ch:implementation}
	%In this chapter we will detail the steps taken to implement the methods discussed in chapter \ref{ch:design}. 
%%%%%%%%%%%%%%%%%%%%%%%%%%%%%%%%%%%%
	\section{Constraining Grid Dimensions}
	\label{sec:grid_constraints}
%%%%%%%%%%%%%%%%%%%%%%%%%%%%%%%%%%%%
There are two constraints that the grid dimensions must conform to. The first is set by the requirements of a simple z-order curve, the second is set by the size of the on chip shared memory. We define $n_{virtual}$ the ratio of a grid dimension and a bin dimension. These constraints are expressed mathematically through the grid dimensions, $n_r, n_{\theta}, n_{\psi}$, and the bin sub-domain dimensions, $nb_r, nb_{\theta}, nb_{\psi}$. As shown in equation \ref{eqn:constraints}, the ratio in each direction must be equivalent and be a power of 2.
		
\begin{equation}
\frac{n_r}{nb_r} = \frac{n_{\theta}}{nb_{\theta}} = \frac{n_{\psi}}{nb_{\psi}} = n_{virtual}
\label{eqn:constraints}
\end{equation} 



The second constraint on the grid dimensions is set by the hardware. The goal is to maximize the shared-multiprocessor occupancy for the charge assign stage of the code. Given that each thread block has the maximum number of threads, 512, and each thread requires roughly 25 registers, the maximum number of thread blocks that can exist simultaneously on a single SM is 2. This means that each thread block can be allocated half of the total amount of shared memory on the SM. Compute capability 2.0 GPUs have 49152 bytes of shared memory per SM. Running two thread blocks per SM provides each block with 24576 bytes of shared memory each, or 6144 floats per block. The maximum that all three bin dimensions can be is 18. For the sake of simplicity these hardware constraints set $nb_r, nb_{\theta}$, and $nb_{\psi} \le 18$. For now none of the bin dimensions are allowed to exceed 18, even if another is smaller. 

A third, loose constraint can be set in order to force a minimum number of thread blocks for the charge assign. The command line option ``--minbins\#'' sets the parameter $n_{virtual} = \#$. This is useful in ensuring that enough thread blocks are launched to populate all of the SMs on the GPU. To populate all of the SMs on a GTX 470 the code would need to launch at least 28 thread blocks. For a GTX 580 with 16 SMs 32 thread blocks are required to fill all of the processors.  
		%\section{Holding to Grid Constraints}

	\section{Particle List Transpose}
As previously mentioned the particle list structure on the GPU is different than the structure on the CPU. On the GPU particles are stored in a structure of arrays, while on the CPU they are stored in a 6x$n$ array. This means that in order to copy a particle list generated on the CPU to the GPU, or vice versa, the particle list must be transposed. The two main places in the code where this matters is when the particle list is initially populated at the start of the code, and when copying a list of pre-calculated reinjection particles from the CPU to the GPU at every time step during the advancing phase.

The particle list transpose was implemented on the CPU in two different ways depending on the compiler used and the available libraries. A GPU based particle list transpose is significantly faster than a CPU based transpose. However, the GPU has a very limited amount of DRAM compared to the CPU, and it is preferable to use as much of the available GPU memory as possible for the main particle list. In any case transposing the entire particle list only occurs once, but a smaller transpose is performed every time step for reinjected particles. This means that while a faster transpose is preferable, it represents so little of the total computation time that it is not worth developing a complicated in place GPU transpose.  

	\section{Charge Assign}
	\label{sec:charge_assign}
	As previously mentioned, the charge assign is one of the most difficult functions to parallelize. The naive approach of applying a thread to every particle and atomically adding each particles contribution to an array in global memory is very slow. Grouping the particles spatially allows the majority of the atomic operations to be done in the context of shared memory which is much faster than global memory. The resulting algorithm resembles basic domain decomposition where each thread block represents a separate sub-domain. The actual charge deposition method in this scheme is very similar to the naive approach, with a key difference being that all the threads in the thread block are operating on shared memory. Once all particles in the sub-domain have contributed to the grid in shared memory it takes only a small number of global memory accesses to write the contributions of a large number of particles to the main array. 

	\subsection{Domain Decomposition}
		The primary grid is decomposed into sub-domains of size $nb_r, nb_{\theta}, nb_{\psi}$. The methods for determining the size of the sub-domains is outlined in section \ref{sec:grid_constraints}. The indexing of the sub-domains is done using a z-order curve in order to preserve spatial locality of the sub-domains in memory. This is done in an attempt to reduce the mean distance that particles must be moved in memory during the sort phase. A graphical representation of this is shown in figure \ref{fig:domain_decomp}.

\begin{figure}
\begin{center}
\includegraphics[width=5in]{implementation/zorder_sceptic.pdf}
\end{center}
\caption[ParticleBin Organization]{Graphical Representation of domain decomposition and ParticleBin organization. The domain is divided up such that every region, highlighted by alternating blue and green fill, is the same size. Additionally the number of bins in any given direction is a power of two. This figure shows what $4^3$ bins looks like.}
\label{fig:domain_decomp}
\end{figure}

In addition to representing a sub-section of the computational mesh, each sub-domain must have a section of the particle list associated with it. The sub-domain must know all of the particles that reside within the region defined by that sub-domain. In essence each sub-domain represents a bin of particles that corresponds to some spatial location, hence the use of ``ParticleBin"  as the naming convention for these object.



	\subsection{Particle Bins}
The \emph{ParticleBin} object keeps track of all of the particles that reside in the region of space that the \emph{ParticleBin} represents. For the sake of simplicity all of the particle bins are the same size spatially, which means that the \emph{ParticleBin} object only has to keep track of the section of the main particle list that the bin represents and the spatial origin of the bin. 

In the context of the particle list each bin represents a pair of bookmarks that bound a section of the particle list. The bookmarks for each bin are calculated after the particle list is sorted by algorithm \ref{alg:count_bins}.

\begin{algorithm}
	\caption{ParticleBin Bookmark Calculation}
	\label{alg:count_bins}
	\begin{algorithmic}
		\FORALL{  $\mathrm{threadID} = 0 \to \mathrm{ParticleList.nptcls}$ in \underline{parallel}}
			\STATE 
			\STATE $\mathrm{binID} = \mathrm{ParticleList.binID}[\mathrm{threadID}]$
			\STATE $\mathrm{binID_{left}} = \mathrm{ParticleList.binID}[\mathrm{threadID} - 1]$
			\STATE $\mathrm{binID_{right}} = \mathrm{ParticleList.binID}[\mathrm{threadID} + 1]$
			
			\IF{$\mathrm{binID} \neq \mathrm{binID_{left}}$}
				\STATE $\mathrm{ParticleBins[binID].ifirstp} = \mathrm{threadID}$
				\STATE $\mathrm{ParticleBins[binID_{left}].ilastp} = \mathrm{threadID} - 1$
			\ENDIF

			\IF{$\mathrm{binID} \neq \mathrm{binID_{right}}$}
				\STATE $\mathrm{ParticleBins[binID].ilastp} = \mathrm{threadID}$
				\STATE $\mathrm{ParticleBins[binID_{right}].ifirstp} = \mathrm{threadID} + 1$	
			\ENDIF
			
		\ENDFOR
	\end{algorithmic}
\end{algorithm}


The spatial origin of the bin is hashed using a z-order curve and stored as a 16-bit unsigned-integer. This 16-bit integer is refered to as the binID and is used for determining the region of the domain that a bin is responsible for as well as a sorting key for the particle list. Calculating the binID will be discussed in more detail in section \ref{sec:plist_sort}. A 16-bit unsigned integer is used for several reasons. First the sorting method detailed in section \ref{sec:plist_sort} is dependent on the number of bits of the sorting key. Second, the upper bound on the grid size set by using a 16-bit integer to store the z-order hash is much larger than the largest grid size that would need to be run. For a 16-bit integer this upper bound is $512^3$ grid points. The third reason for using a 16-bit integer is that it also reduces the memory requirements of the particle list by about 5\%, which does help when trying to run as many particles on the GPU as possible. 

	\subsection{Particle Push}
Now that the particles are organized spatially in memory, it is trivial to assign a single thread block to a region of space and corresponding particle bin in order to perform the particle push. This process is rather simple and is outlined in psuedo code in algorithm \ref{alg:gpu_c2mesh}. 

\begin{algorithm}
	\caption{GPU Charge Assign}
	\label{alg:gpu_c2mesh}
	\begin{algorithmic}
		\FORALL{ $ParticleBin \in Grid$ in \underline{parallel}}
			\STATE \emph{$\backslash \backslash$ Inside the threadBlock with ID blockID}

			\STATE \_\_shared\_\_ float $subGrid(nb_r,nb_{\theta},nb_{\psi})$

			\FORALL{ $node \in subGrid$ in \underline{parallel}}
			\STATE $node = 0$
			\ENDFOR
			
			\STATE \_\_syncthreads()
			

			\FORALL{  $particle \in ParticleBin$ in \underline{parallel}}
				\STATE cell = $particle.cell - ParticleBin.origin$
				\FORALL{ $node \in cell$}
					\STATE atomicAdd($subGrid(cell,node)$, weight($node$))
				\ENDFOR
			\ENDFOR
			\STATE \_\_syncthreads()

			\STATE \emph{$\backslash \backslash$ Write block results to global memory}
			\FORALL{ $node \in subGrid$ in \underline{parallel}}
				\STATE atomicAdd($Grid(blockID,node)$, $subGrid(node)$)
			\ENDFOR
			
		\ENDFOR
	\end{algorithmic}
\end{algorithm}

Each thread block reads in 512 particles at a time, although only 32 particles, a warp, are processed in parallel within the block. Each thread within this warp loops over the 8 nodes that bound the cell that contains the particle being processed. The nodes reside in a shared memory array, and are updated with the weighted particle data atomically. Once all of the nodes for a given particle have been updated the thread will retrieve a new particle from global memory. This process is repeated by all of the threads in the block until every particle in the particle bin has been processed. Once all of the particles have been processed the block then atomically updates the nodes in global memory with the values stored in shared memory.  

The atomic operations in this algorithm lead to some very interesting time complexity behavior. In essence this algorithm is being executed on a machine with 32 processors. The time complexity of this scenario is $\mathcal{O}(\frac{c}{p})$, where $c$ is constant and $p$ is the number of available processors. When two processors attempt to atomically update the same memory address, one of the processors must wait until the other is finished. This means that one processor is effectively lost for a 1-way conflict. 

The mean number of n-way atomic conflicts $N$ in a warp over a sub domain of size $G$, and the execution time $T$ is given by:
\begin{equation}
N = \frac{31!}{(31-n)! G^n} \hspace{0.1\textwidth} T(n) \propto \frac{c}{32-n}
\end{equation}

This means that the total time complexity of this algorithm with respect to the sub domain size $G$ is:

\begin{equation}
T(G) \propto c \cdot \sum_{n=1}^{31} \frac{1}{32-n} \frac{31!}{(31-n)! G^n}
\end{equation}

This behavior can be seen roughly in \ref{fig:subdomain_size_scan}. This algorithm on the GPU can perform the particle push up to 200x faster than the CPU version of the charge assign. However, this algorithm relies on the particle data being ordered spatially, which contributes to the run time. The method used to maintain an ordered particle list on the gpu will be discussed in the following section.


	\section{Particle List Sort}
	\label{sec:plist_sort}
	
	As previously mentioned in section \ref{sec:charge_assign} an ordered particle list must be maintained in order for the charge assign to be fast. The particle list sort, algorithm \ref{alg:sort_overview} consists of three distinct subroutines, populating the key/value pairs, sorting the key/value pairs, and finally a payload move. 

\begin{algorithm}
	\caption{Particle List Sort Overview}
	\label{alg:sort_overview}
	\begin{algorithmic}
		\STATE
		\STATE Populate\_KeyValues(Particles, Mesh, sort\_keys, sort\_values)
		\STATE
		\STATE thrust::sort\_by\_key(sort\_keys,sort\_keys+nptcls,sort\_values)
		\STATE
		\STATE Payload\_Move(Particles, sort\_values)
	\end{algorithmic}
\end{algorithm}

This method of maintaining particle list order was chosen because it is a good balance between simplicity and performance. An additional benefit of this routine is that it uses the sort from the \gls{gls:thrust} library, which is maintained by NVIDIA. 

		\subsection{Populating Key/Value Pairs}
The first step in sorting the particle list is ensuring that the key/value pairs needed by the sorting routine are populated. The sorting key for a particle is the index of the particle bin that the particle belongs to. Sorting values are simply the position of the particle in the unsorted list. 

Calculating the particle bin index, or binid, begins with calculating the mesh cell that the particle resides in. This cell described by coordinates $i_r, i_{\theta}, i_{\phi}$. The coordinates of the particle bin that a given cell resides in is given by:

\begin{equation}
ib_r = \frac{i_r}{nb_r} \mathrm{;}\hspace{0.5in} ib_{\theta} = \frac{i_{\theta}}{nb_{\theta}} \mathrm{;} \hspace{0.5in} ib_{\phi} = \frac{i_{\phi}}{nb_{\phi}}
\end{equation}

	The resulting block coordinates are then hashed using a z-order curve and stored as the binid. Each thread calculates the binid's for several particles and stores them in the sort\_keys array. Once a thread has calculated the binid for a particle it also stores the index of that particle as an integer in the sort\_values array. 

		\subsection{Sorting Key/Value Pairs}
		The key/value pair sorting is done using the \gls{gls:thrust} library sort\_by\_key template function. This function is provided by NVIDIA with \gls{gls:CUDA}. The \gls{gls:thrust} sort is a radix sort that has been optimized for NVIDIA GPUs\cite{NVIDIACorporation2011a}. The snippet of the sort code used in SCEPTIC3DGPU is shown in figure \ref{fig:thrust_sort}.

\begin{figure}
\begin{lstlisting}[frame=single]
	// wrap raw device pointers with a device_ptr
	thrust::device_ptr<ushort> thrust_keys(binid);
	thrust::device_ptr<int> thrust_values(particle_id);

	// Sort the data
	thrust::sort_by_key(thrust_keys,thrust_keys+nptcls,thrust_values);
	cudaDeviceSynchronize();
\end{lstlisting}
\vspace{-0.4in}
\caption{Thrust Sort Setup and Call}
\label{fig:thrust_sort}
\end{figure}


		\subsection{Payload Move}
		The payload move is responsible for moving all of the particles from their old locations in memory to the new sorted locations. The idea is simple, each thread represents a slot on the sorted particle list. Threads read in an integer, the particleID, from the values array that was sorted using the binid's. This integer is the location of a given threads particle data in the unsorted list. Data at index particleID is read in, and stored in the new list at index threadID. While the idea is simple, this algorithm would require a completely separate copy of the particle list, a lot of wasted memory. However, since the particle list is set up as a structure of arrays, there is something that can be done to significantly reduce the memory requirements. The method, outlined in algorithm \ref{alg:payload_move} reorders only a single element of the particle list structure at a time. 

\begin{algorithm}
	\caption{GPU Payload Move}
	\label{alg:payload_move}
	\begin{algorithmic}
		\FORALL{ $member \in XPlist$}
		\STATE float* idata = member
		\STATE float* odata = XPlist.buffer
		\STATE \textbf{reorder\_data}(odata,idata,particleIDs)
		\STATE member = odata
		\STATE buffer = idata
		\ENDFOR
	\end{algorithmic}
\end{algorithm}
  		
		Essentially the idea is that a great deal of memory can be saved by statically allocating a ``buffer'' array that is the same size as each of the data arrays. During the payload move each data array is sorted into the buffer array. Some pointer swapping is performed, the old buffer array becomes the new data array, and the old data array becomes the buffer for the next data array. For SCEPTIC3DGPU this implementation of the payload move only increases the particle list size by about 8.6\%.



	\section{Poisson Solve}
The field solve used in SCEPTIC3D involves solving the electrostatic Poisson equation:

\begin{equation}
\nabla^2\phi=\frac{n_e-n_i}{\Lambda_{De}^2}
\label{eqn:poisson_eqn1}
\end{equation}

Where $\phi$ is the electrostatic potential, $n_{e,i}$ the electron, ion charge densities normalized to the background charge density $N_\infty$, and $\Lambda_{De}$ is the unperturbed electron Debye length $\Lambda_{De} = \sqrt{\varepsilon_0T_e/N_\infty e^2}$. Using the SCEPTIC3D approximation that the electron density can be taken to be a Boltzmann distribution, and normalizing $\nabla$ to $1/R_p$ the probe radius, the Poisson equation can be rewritten as 

\begin{equation}
\nabla^2\phi=\frac{exp(\phi)-n}{\lambda_{De}^2}
\label{eqn:poisson_eqn2}
\end{equation}

We are changing $n_i$ to $n$ in order to avoid confusion with the use of $i$ as a spatial index. Assuming that the potential varies little from one time step to the next we can expand equation \ref{eqn:poisson_eqn2} about the known potential $\phi^*$ yielding
   
\begin{equation}
\nabla^2\phi_{i+1}=\frac{exp(\phi^*)[1+(\phi_{i+1}-\phi^*)]-n}{\lambda_{De}^2}
\label{eqn:poisson_eqn3}
\end{equation}

Equation \ref{eqn:poisson_eqn3} can be linearized into the form $A\phi+\omega = \sigma$ using finite volumes. Following the derivation for the linear operator $A$ from \cite{Patacchini2010} we get:


\begin{align}
(A\phi)_{i,j,k} &= a_i\phi_{i+1,j,k}
		  -b_i\phi_{i-1,j,k}
		  +c_{i,j}\phi{i,j+1,k}
		  -d_{i,j}\phi_{i,j-1,k} \nonumber \\
		  &\qquad{} +e_{i,j}(\phi_{i,j,k+1}-\phi_{i,j,k-1})
		  -[f_{i,j} + exp(\phi_{i,j,k}^*)]\phi_{i,j,k}
\label{eqn:psolve}
\end{align}

The full derivation of equation \ref{eqn:psolve} along with the definitions of the vectors $a$ through $f$ can be found in \cite{Haakonsen2011}.

Now that we have our linear operator defined, along with the vectors $\omega$ and $\sigma$ we can see that this is a simple sparse linear system. These system is fairly easy to solve using any sparse linear system solving scheme. One of the easier schemes to implement on the GPU is the preconditioned biconjugate gradient method \emph{PBCG}. This method is easy to parallelize due to the simplicity of the operations involved, namely \emph{PBCG} consists solely of sparse matrix-vector products, vector dot products, and multiplication of vectors by a scalar. All of these operations can be implemented efficiently on the GPU. 

We can make an additional optimization of the PBCG solver by recognizing that only one element of our linear operator $A$ changes every time step. This means that we can make use of some of the read-only memory spaces on the GPU, such as texture memory. Since texture memory is cached spatially in its own on-chip cache, we hope that we can significantly reduce the number of global memory transactions for all of the matrix multiplications. We ended up implementing two versions of the Poisson solve, one that stores both $A$ and $\phi^*$ as textures, and one that stores them in global memory. The version to be used is determined at compile time with the use of pre-processor macros. 
The only concern with the texture implementation would be the cost of calling cudaBindTextureToArray(). However, since only one diagonal of $A$ and and $\phi^*$ change over the course of the time step, the binding cost should not be an issue. Using textures for arrays in other parts of the code is also possible. Since we are using textures for the potential in the field solve it is very simple to also use a texture based potential for the particle advancing. The actual gain of using textures will be discussed in chapter \ref{ch:performance}.

We will not develop GPU implementations of the pre-processing steps required to prepare the ion density $n$ from the particle-density interpolation since those steps do not represent a significant computational cost. Additionally, the particle-density array already exists in host memory due to the fact that an \gls{ac:mpi} reduce must be performed to gather this array from multiple GPUs. Since this data is already in host memory, and the computational time is low, there is very little reason to rewrite these routines for the GPU.  


 
%%%%%%%%%%%%%%%%%%%%%%%%%%%%%%%%%%%%
	\section{Particle List Advance}
%%%%%%%%%%%%%%%%%%%%%%%%%%%%%%%%%%%%
	
Moving the particles on the grid starts with determining the acceleration of the particle. This is calculated by interpolating the gradient of the potential, $\nabla\phi$, from the spherical mesh using the same methods as the cpu code. The new position of the particle is then calculated using the leap-frog method:

\begin{equation}
\begin{aligned}
\mathbf{v}' &= \mathbf{v}+\mathbf{a}\Delta t \\
\mathbf{x}' &= \mathbf{x}+\mathbf{v}'\Delta t
\end{aligned}
\end{equation}
 Additional details of the particle advance can be found in reference \cite{Patacchini2007} section 3.1.2.
		
While the implementation of the basic physics of the particle advance remains the same, there were several issues. Quickly determining whether a particle has crossed one of the domain boundaries, contributing to diagnostic outputs, and handling reinjections were the main issues. 

		\subsection{Checking Domain Boundaries}
In order to correctly contribute to the diagnostic outputs, the location where a particle left the domain must be known. This means that the process of checking whether or not a particle has left the grid must also calculate the position of the particle when it crossed the boundary. Since the boundaries are both spheres, we can apply a very common ray tracing technique used to calculate a line-sphere intersection. 

The equation for the path of a particle with initial position $\mathbf{p}_i$ and final position $\mathbf{p}_{i+1}$ is given by;

\begin{equation}
 \mathbf{l} = \frac{\mathbf{p}_{i+1} - \mathbf{p}_{i}}{\|\mathbf{p}_{i+1} - \mathbf{p}_{i}\|}
\mathbf{x} = d\mathbf{l}
\label{eqn:p_direction}
\end{equation}

Combining equation \ref{eqn:p_direction} to the equation of a sphere with radius $r$ centered at the origin, and expanding we get:

\begin{equation}
\begin{aligned}
\|d\mathbf{l} - (-\mathbf{p}_i)) \|^2  &= r^2 \\
d^2\mathbf{l}^2 + 2d(\mathbf{l}\cdot \mathbf{p}_i) + \mathbf{p}_i^2 &= r^2
		\end{aligned}
\end{equation}

We can then use the quadratic equation to solve for $d$. There are a few qualifiers, namely, $(\mathbf{l}\cdot \mathbf{p}_i)^2 - \mathbf{p}_i^2 + r^2 \geq 0$. The distance to the nearest intersection is simply the smallest positive solution of $d$. A particle leaves the domain if $d^2 \le \|\mathbf{p}_{i+1} - \mathbf{p}_{i}\|^2$. If this condition is true, then the particle's new position $p_{i+1}$ is changed such that $d^2 = \|\mathbf{p}_{i+1} - \mathbf{p}_{i}\|^2$. The particle's position is changed in order to accommodate some of the diagnostic outputs discussed in section \ref{sec:diagnostic_outputs}. 



		\subsection{Handling Reinjections}
Once it has been determined that particles have left the grid, new particles must be reinjected to replace them. In the serial version of the code this is handled by simply calling a reinjection subroutine that determines the new particle's position and velocity. Once the new position and velocity has been found, the particle is moved for the remainder of the time step, and replaced by a new particle if the reinjected particle leaves the domain.  

Performing reinjections in this manner on the GPU would introduce very large divergences in warp execution as well as very uncoalesced memory accesses. Eliminating the warp divergences would require that all of the threads in a warp be operating on reinjected particles. Reducing the uncoalesced memory accesses would require that all of the reinjected particles be adjacent in memory.  Since we already have an object with methods that can move a list of particles and handle reinjections, all we really need is some method by which we can efficiently and reversibly ``pop" a subset of the particles in the main list to a secondary list. From there we can perform the particle advance on the secondary list, and place the results back in the empty particle slots in the main list. The resulting advancing algorithm is as follows:

\begin{algorithm}
	\caption{Particle Advancing Algorithm}
	\label{alg:advancing}
	\begin{algorithmic}
		\STATE // Update The particle positions and check domain boundaries
		\STATE $\textbf{GPU\_Advance}(particles,mesh,Exit\_Flags)$
		\STATE
		\STATE \textbf{Prefix\_Scan}(Exit\_Flags)
		\STATE
		\STATE nptcls\_reinject = Exit\_Flags[nptcls-1]
		\STATE
		\IF{nptcls\_reinject $>$ 0}
			\STATE
			\STATE reinjected\_particles.allocate(nptcls\_reinject)
			\STATE
			\STATE \textbf{Stream\_Compact}(particles $\subset$ exited $\rightarrow$ reinjected\_particles)
			\STATE
			\STATE // Recursively call this algorithm on the reinjected particles
			\STATE reinjected\_particles.advance()
			\STATE
			\STATE \textbf{Stream\_Expand}(reinjected\_particles $\rightarrow$ exited $\supset$ particles)
		\ENDIF
		\STATE
		\STATE return
	\end{algorithmic}
\end{algorithm}

A graphical representation of algorithm \ref{alg:advancing} can be seen in figure \ref{fig:padvnc}, which also includes the handling of collisions. For the handling of both collisions and reinjections four separate steps are needed. Step (1) determines whether or not a particle will be colliding, when the collision happens, and flags that particle in the list. Step (2) advances the particle's positions and velocities, and flags all particles that leave the simulation domain. Particles undergoing collisions in step (2) are only advanced to the time of collision, and particles that leave the domain are advanced to the point where they exit the domain. Particles that are flagged for a collision, but leave the domain before the collision takes place do not undergo a collision, but are flagged for reinjection. At this point we have several groups of particles, two that require special operations and additional advancing, and a third that is completely finished. 
In step (3) we separate these groups of particles and into compacted collision and reinjection lists, and a sparse main list. The compacted lists are operated by the collision and reinjection operators respectively and then moved recursively. 
Once the recursive moves have been completed the particles in the two sub-lists are placed back in their original places in the main list by step (4). It should be noted that collisions have not yet been implemented in the GPU code, but the data handling infrastructure needed to support them as shown in figure \ref{fig:padvnc} has been implemented. 
Since the collision and reinjection operators are completely independent, we could make use of multiple \gls{gls:CUDA} ``streams" such that the two processes could be carried out asynchronously on the same GPU. 

\begin{figure}
\begin{center}
% Turing Machine
% Author: Sebastian Sardina
%\documentclass[a4paper,10pt]{article}
%\usepackage[usenames,dvipsnames]{xcolor}
%\usepackage{tikz}
%\usepackage{ifthen}
%\usetikzlibrary{chains,fit,shapes}

%\begin{document}


\noindent
\begin{tikzpicture}
		\tikzstyle{every path}=[very thick]
		\tikzstyle{MainSty}=[->,dashdotted]
		\tikzstyle{CollSty}=[->,color=Emerald,densely dashdotted]
		\tikzstyle{ReinSty}=[->,color=Maroon,dashed]

		\edef\sizetape{1.0cm}
		\tikzstyle{tmtape}=[draw,minimum size=\sizetape]
		\tikzstyle{tmtape1}=[draw,minimum height=\sizetape,minimum width=0.5cm]
		\tikzstyle{tmhead}=[arrow box,draw,minimum size=0.4cm,arrow box
		arrows={east:.25cm}]

		\tikzstyle{cfill1}=[fill=Purple]
		\tikzstyle{cfill2}=[fill=SkyBlue]
		\tikzstyle{cfill3}=[fill=Maroon]
		\tikzstyle{cfill4}=[fill=Emerald]


\def\MainArray{10,9,8,7,6,5,4,3,2,1}
\def\ReinArray{8,5,1}
\def\CollArray{9,5,4,2}

\newcounter{prim@a}
\newcounter{prim@b}
\newcounter{prim@c}
\newcounter{prim@d}


% Figure out which nodes will collide

\begin{scope}[start chain=5,node distance=1cm]

% The initial list
	%\node [on chain=5]	{\begin{tikzpicture}[remember picture]
		\begin{scope}[start chain=1 going below,node distance=-1pt]

		\foreach \i in \MainArray {
			\node [on chain=1,tmtape] (pmain\i) {$\i$};
		}
		\end{scope}
		%\end{tikzpicture}};

% The Collision Check
	%\node [on chain=5]	{\begin{tikzpicture}[remember picture]
		\begin{scope}[xshift=2.0cm,start chain=1 going below,node distance=-1pt]


		\foreach \i in \MainArray{
			
			\foreach \j in \CollArray
			{
				\setcounter{prim@a}{\i}
				\setcounter{prim@b}{\j}
				\ifthenelse{\value{prim@a}=\value{prim@b}}
				{\setcounter{prim@c}{\i}}{}%
			}

			\ifthenelse{\value{prim@a}=\value{prim@c}}
			{ \node [on chain=1,tmtape1,fill=Emerald] (pc\i) {};}%
   		{ \node [on chain=1,tmtape1] (pc\i) {};}%

		}
		\end{scope}
		%\end{tikzpicture}};

% The Reinjection Check
	%\node [on chain=5]	{\begin{tikzpicture}[remember picture]
		\begin{scope}[xshift=3.75cm,start chain=1 going below,node distance=-1pt]

		\foreach \i in \MainArray{

			\foreach \j in \ReinArray
			{
				\setcounter{prim@a}{\i}
				\setcounter{prim@b}{\j}
				\ifthenelse{\value{prim@a}=\value{prim@b}}
				{\setcounter{prim@c}{\i}}{}%
			}

			\foreach \j in \CollArray
			{
				\setcounter{prim@a}{\i}
				\setcounter{prim@b}{\j}
				\ifthenelse{\value{prim@a}=\value{prim@b}}
				{\setcounter{prim@d}{\i}}{}%
			}

			\ifthenelse{\value{prim@a}=\value{prim@c}}
			{ \node [on chain=1,tmtape1,fill=Maroon] (pr\i) {};}%
   		{ \ifthenelse{\value{prim@a}=\value{prim@d}}
				{\node [on chain=1,tmtape1,fill=Emerald] (pr\i) {};}
				{\node [on chain=1,tmtape1] (pr\i) {};}

			}%

		}
		\end{scope}
		%\end{tikzpicture}};

% The initial list
	%\node [on chain=5]	{\begin{tikzpicture}[remember picture,start chain=4 going above,node distance=0.5cm]
	\begin{scope}[xshift=7.5cm,yshift=-9.0cm,start chain=4 going above,node distance=0.5cm]
		\node [on chain=4,tmtape] (wait) {Wait};

		\begin{scope}[yshift=9.0cm,start chain=1 going below,node distance=-1pt]
			\node [on chain=1] {\footnotesize Collide \& Move};

			\foreach \i in \CollArray {

			\foreach \j in \ReinArray
			{
				\setcounter{prim@a}{\i}
				\setcounter{prim@b}{\j}
				\ifthenelse{\value{prim@a}=\value{prim@b}}
				{\setcounter{prim@c}{\i}}{}%
			}

				\ifthenelse{\value{prim@a}=\value{prim@c}}{}
				{\node [on chain=1,tmtape] (lc\i) {$\i$};}%
			}
		\end{scope}

		\begin{scope}[yshift=5.0cm,start chain=1 going below,node distance=-1pt]

			\node [on chain=1] {\footnotesize Reinject \& Move};
			\foreach \i in \ReinArray {
				\node [on chain=1,tmtape] (lr\i) {$\i$};
			}
			\end{scope}

		\end{scope}
		%\end{tikzpicture}};

% The Main list without collision and reinjection
	%\node [on chain=5]	{\begin{tikzpicture}[remember picture]
		\begin{scope}[xshift=11.325cm,start chain=1 going below,node distance=-1pt]

		\foreach \i in \MainArray{

			\foreach \j in \ReinArray
			{
				\setcounter{prim@a}{\i}
				\setcounter{prim@b}{\j}
				\ifthenelse{\value{prim@a}=\value{prim@b}}
				{\setcounter{prim@c}{\i}}{}%
			}

			\foreach \j in \CollArray
			{
				\setcounter{prim@a}{\i}
				\setcounter{prim@b}{\j}
				\ifthenelse{\value{prim@a}=\value{prim@b}}
				{\setcounter{prim@c}{\i}}{}%
			}

			\ifthenelse{\value{prim@a}=\value{prim@c}}
			{ \node [on chain=1,tmtape] (pm\i) {$\i^*$};}%
   		{ \node [on chain=1,tmtape] (pm\i) {$\i$};}%

		}
		\end{scope}
		%\end{tikzpicture}};


		\foreach \i in \MainArray
		{

			\setcounter{prim@a}{\i}

			\foreach \j in \ReinArray
			{
				\setcounter{prim@a}{\i}
				\setcounter{prim@b}{\j}
				\ifthenelse{\value{prim@a}=\value{prim@b}}
				{\setcounter{prim@c}{\i}}{}%
			}

			\foreach \j in \CollArray
			{
				\setcounter{prim@b}{\j}
				\ifthenelse{\value{prim@a}=\value{prim@b}}
				{\setcounter{prim@d}{\i}}{}%
			}

			\ifthenelse{\value{prim@a}=\value{prim@c}}
			{ 
				\draw[MainSty] (pmain\i) to[->] (pc\i);

			  \ifthenelse{\value{prim@a}=\value{prim@d}}
				{\draw[CollSty] (pc\i) to (pr\i);}
				{\draw[MainSty] (pc\i) to (pr\i);}
			  \begin{scope}[ReinSty]
				\draw[->] (pr\i.mid east) to[in=180] (lr\i);
				\draw[->] (lr\i) to[out=0,in=180] (pm\i.mid west);
			  \end{scope}

			}%
   		{\ifthenelse{\value{prim@a}=\value{prim@d}}
			{ \draw[MainSty] (pmain\i) to (pc\i);
			  \begin{scope}[CollSty]
				\draw[->] (pc\i) to (pr\i);
				\draw[->] (pr\i.mid east) to[in=180] (lc\i);
				\draw[->] (lc\i) to[out=0] (pm\i.mid west);
			  \end{scope}
			}
			{	\begin{scope}[MainSty]
				\draw[->] (pmain\i) to[->] (pc\i);
				\draw[->] (pc\i) to (pr\i);
				\draw[->] (pr\i.mid east) to[in=180] (wait);
				\draw[->] (wait) to[out=0] (pm\i.mid west);
				\end{scope}
			}%
			}

		}

		\begin{scope}[xshift=2.0cm,yshift=1cm,start chain=1,node distance=0.875cm,align=center]

			\node [on chain=1,tmtape,draw=none] {\large\textbf{(1)}};
			\node [on chain=1,tmtape,draw=none] {\large\textbf{(2)}};
			\node [on chain=1,tmtape,draw=none] {};
			\node [on chain=1,tmtape,draw=none] {\large\textbf{(3)}};
			\node [on chain=1,tmtape,draw=none] {};
			\node [on chain=1,tmtape,draw=none] {\large\textbf{(4)}};
	
		\end{scope}







\end{scope}

\end{tikzpicture}
%\end{document}


\end{center}
\caption[Illustration of GPU particle advancing algorithm]{Illustration of GPU particle advancing algorithm with handling of reinjections and collisions. All of the particles to be advanced begin in the same list. Prior to updating the particle's positions and velocities, step (1) determines which particles will undergo collisions, and when those collisions will occur. Between steps (1) and (2) the particle's positions and velocities are advanced. After the advance some particles will have left the domain, these particles are flagged in step (2), with exit flags overriding collision flags. After step (2) two stream compactions are performed in order to condense the re-injection and collision lists. In step (3) the collision and reinjection operators are applied to these two lists respectively, and the entire move method is called recursively on each of the lists. Once all recursive steps have finished the collision and re-injection lists are merged back into the main array in step (4). }
\label{fig:padvnc}
\end{figure}

Compacting some subset of a parent list is a fairly easy parallel operation called stream compaction. Stream compaction, shown in figure \ref{fig:stream_compact} is a process by which a random subset of a list can be quickly copied to a new list in parallel. It only works for some binary condition, such as an array of length nptcls, where each element is 1 for particles that have left the domain, and 0 for all others. For each `true' element taking the cumulative sum of all preceding elements yields a unique number that can be used as an index in a new array. 

\begin{figure}
\begin{center}
\includegraphics[width=4in]{implementation/stream_compact.jpg}
\end{center}
\caption[Stream Compaction]{Stream Compaction from GPU Gems 3 \cite{Harris2007}}
\label{fig:stream_compact}
\end{figure}



The new positions and velocities for reinjected particles are taken from a precalculate pool of new particles. This pool is approximately 1/10\textsuperscript{th} the size of the main particle list, and is repopulated prior to every advance step. Prior to the first advance step the entire pool is populated, but subsequent steps only refill slots that have been used in reinjections. This ``pool" method was chosen over implementing the reinjection calculations on the GPU for reasons of code maintenance. There are currently several different reinjection schemes that SCEPTIC3D can call. Maintaining versions of those routines in both fortran and \gls{gls:CUDA} would be more cumbersome than the small benefit that calculating them on the GPU would provide. 

		\subsection{Diagnostic Outputs}
		\label{sec:diagnostic_outputs}
	There are several outputs from the particle advancing routine that are desirable to keep in the GPU code. These outputs include;


\begin{itemize}
\singlespacing
\item Total momentum transferred to the probe.
\item Total momentum lost to particles leaving the system.
\item Total momentum added to the system from reinjections.
\item Total energy transferred to the probe. 
\item Total number of particles lost to the probe. 
\item Spatial distribution of energy transferred to the probe. 
\item Spatial distribution of momentum transferred to the probe. 
\item Spatial distribution of density transferred to the probe. 
\end{itemize}


	For the single value quantities, such as the quantities of total momentum, energy, and particle count are tallied using atomic operations in shared memory. Since the fraction of particles that leave the domain during a timestep is small, $\leq 7\%$, we can just use conditional statements within the primary move kernel to perform the tallies. While this is not strictly optimal, it does not have a very large impact on the overall performance of the move kernel. 







\chapter{Performance} 

	\section{Test Setup}
		\subsection{Parameter Space Explored}
		\subsection{Machine Parameters}
		\subsection{Memory Bandwidth Comparison}
	
	\section{Particle list size scan}
The following tests were performed to explore the dependence of sceptic3Dgpu's runtime on the number of particles per node. These tests were performed using two MPI nodes and a grid size of 128x64x64.

\begin{figure}
\begin{center}
\includegraphics[width=6in]{performance/nptclsize_scan128x64x64ons8bins.pdf}
\end{center}
\caption{Number of Particles Scan on a 128x64x64 grid}
\label{fig:nptclsize_scan128x64x64}
\end{figure}

\begin{figure}
\begin{center}
\includegraphics[width=6in]{performance/nptclsize_scan64x32x32ons8bins.pdf}
\end{center}
\caption{Number of Particles Scan on a 64x32x32 grid}
\label{fig:nptclsize_scan128x64x64}
\end{figure}

\begin{figure}
\begin{center}
\includegraphics[width=6in]{performance/nptclspeedup_scan128x64x64ons8bins.pdf}
\end{center}
\caption{Speedup factor Number of Particles Scan on a 128x64x64 grid}
\label{fig:nptclsize_scan128x64x64}
\end{figure}


\begin{figure}
\begin{center}
\includegraphics[width=6in]{performance/gridsize_scan8ptcls8bins.pdf}
\end{center}
\caption{Gridsize Scan with 8 million ptcls}
\label{fig:nptclsize_scan128x64x64}
\end{figure}


\begin{figure}
\begin{center}
\includegraphics[width=6in]{performance/gridsize_scan8ptcls16bins.pdf}
\end{center}
\caption{Gridsize Scan with 8 million ptcls}
\label{fig:nptclsize_scan128x64x64}
\end{figure}


\begin{figure}
\begin{center}
\includegraphics[width=6in]{performance/gridsize_scan16ptcls8bins.pdf}
\end{center}
\caption{Gridsize Scan with 16 million ptcls}
\label{fig:nptclsize_scan128x64x64}
\end{figure}

\begin{figure}
\begin{center}
\includegraphics[width=6in]{performance/gridsize_scan16ptcls16bins.pdf}
\end{center}
\caption{Gridsize Scan with 16 million ptcls}
\label{fig:nptclsize_scan128x64x64}
\end{figure}

\begin{figure}
\begin{center}
\includegraphics[width=6in]{performance/gridsize_scan34ptcls8bins.pdf}
\end{center}
\caption{Gridsize Scan with 34 million ptcls}
\label{fig:nptclsize_scan128x64x64}
\end{figure}

\begin{figure}
\begin{center}
\includegraphics[width=6in]{performance/gridsize_scan34ptcls16bins.pdf}
\end{center}
\caption{Gridsize Scan with 34 million ptcls}
\label{fig:nptclsize_scan128x64x64}
\end{figure}

\begin{figure}
\begin{center}
\includegraphics[width=6in]{performance/gridshape_scan.pdf}
\end{center}
\caption{Sub Domain Size scan}
\label{fig:subdomain_size_scan}
\end{figure}

	
	\section{Grid Size scan}
		\subsection{Absolue Size}
		\subsection{Ratio of Surface Area to Volume}
			\subsubsection{Global Domain Ratio}
			\subsubsection{Threadblock Sub-Domain Ratio}

	\section{Kernel Parameters Scan}

	\section{Discussion}

\chapter{Conclusion}

\section{Review}
Towards the end of chapter \ref{ch:introduction} we stated the primary goals of this project. We set out to develop a multi-GPU version of SCEPTIC3D using very generalized techniques. We came away with concrete conclusions concerning the performance impacts of using a full radix sort. We also explored ways to minimize warp divergence while handling reinjections and collisions using stream compaction and recursion. Lastly, we investigated the benefits of utilizing texture memory for constant arrays that are frequently accessed. In the end, the lessons that we learned from the exploration of each of these topics enabled us implement a GPU PIC code that attained a upwards of 100 times the performance of its CPU based counterpart. 

While the end results were impressive, we learned that a PIC code implemented on the GPU can differ greatly from the same code implemented on a CPU. Over the course of this project we analyzed the methods used by others in the development of GPU PIC implementations. The main issue that was identified by most groups was how to efficiently implement the particle to mesh step on the GPU. The conclusion that we, and many of the papers we reviewed, came to was that the particle data must be sorted in order to ensure that the majority of the accumulation can be performed using fast shared memory. So far nearly every paper that we reviewed developed their own techniques to keep the particle list sorted, with each technique optimized to fit a specific problem. 
While these specialized methods can achieve higher performance, more generalized methods tend to be easier to implement, and if they can be based on external libraries, easier to maintain. 

Some of these past works produced very generalized methods. The particle-mesh integration technique developed by Stantchev et al is a good example of a broadly applicable technique. We aimed to develop more general techniques for other steps of the PIC method. We presented a fast, very general sort based on the the radix sort provided by the THRUST library. While the costs associated with this sort are not negligible, the combined costs of the sort and charge assign are still comparable to the costs of the advancing step and the field solve steps. The other general technique that we developed was a way to minimize warp divergence when dealing with reinjections and collisions. The specific algorithm presented works well when the number of collisions and reinjections are small, but can be applied more generally by changing the conditions for placing a particle in a sub-list. 

The final unique property of the GPU PIC implementation presented here is the use of texture memory as the primary data structure for the potential and the backbone of the Poisson solver. We concluded that texture memory should be used as a data structure for spatially organized data. Of course there are drawbacks to texture memory, primarily that the very strict usage properties can make it difficult to integrate with the rest of the code. 



\section{Implications}
We were able to achieve speedup factors of 160 versus older CPUs and factors of 60 versus newer Intel i7 CPUs. With these speedup factors it is possible to attain cluster level performance on a multi-GPU workstation for SCEPTIC3D. In the tests we performed, we ran 50 million particles per GTX 590. There are commercially available workstations that support up to 8 graphics cards. Populated with 8x GTX 590 graphics cards, one of these workstations would be the equivalent of a 2560 core cluster for the purposes of running SCEPTIC3D. Essentially a GPU enabled SCEPTIC3D can attain cluster level performance on a single workstation. 

While we managed to improve the performance of the particle advancing and the particle-mesh interpolation steps by factors of 122 and 104 respectively, we ran into the issue where the Poisson solve became the dominant cost. As we saw in chapter \ref{ch:performance} when the grid becomes much larger than $64^3$ the cost of the Poisson solve dominates. If we start to consider even larger grid sizes, we eventually run into the limit where we are using more memory for the grid than the particle data. At that point we would have to decompose the domain across multiple GPUs such that each GPU only moved particles on part of the grid. This would greatly complicate the entire system because particles that leave one sub-domain would have to be passed from one GPU to another. Smaller PIC codes should not have to worry too much about this, although methods will have to be developed to handle this efficiently for larger codes.

One of the biggest implications of this work is its effect on the particle sorting step. Previous groups attempted to solve the sorting question using the radix sort implementation provided by the CUDPP library. However, this implementation was unacceptably slow, and combined with the poor memory access patterns of older cards, the library sort was out of the question. Recently the radix sort, now provided by the THRUST library, has been significantly improved. This fact, combined with the improved memory access patterns has resulted in a radix based particle sort method with significantly better performance. The performance of the THRUST sort is currently good enough that there is little reason to develop a specialized sorting technique. 

\section{Future Work}
Within the scope of the GPU implementation of SCPETIC there is at least one additional physics component that has not yet been included due to time constraints, the collision operator. While we have laid out the framework for implementing collisions on the GPU we have not yet implemented the entirety of that framework. Additionally, the re-injection particle list is currently calculated on the CPU, and a portion of this 'pool' is transfered to the GPU during the re-injection phase. Refilling the particle list on the CPU currently costs about 1/3\textsuperscript{rd} the cost of the particle advance. Implementing this on the GPU would significantly reduce the cost of this step, as well as eliminate the particle list transpose and host to device memory copy currently required. Further improvements to SCEPTIC3D could be made by optimizing the advancing kernel, specifically optimizing the force interpolation function to require fewer registers. 

Pertaining to GPU PIC implementations in general, more work needs to be done concerning large grid sizes. As previously stated, two of the major issues facing larger scale GPU PIC implementations are multi-GPUs domain decomposition and optimizing large field solves. Developing asynchronous techniques that perform computations on both the GPU and CPU simultaneously is another high priority. 




\appendix
\chapter{Tables}

\begin{center}
\begin{table}
\begin{tabular}{| p{4.0cm} | p{3.5cm} |}
\hline
Component & Runtime (ms) \\ \hline
Particle data read, move, and write & 375 \\ \hline
Potential Grid Read & 467 \\ \hline
Charge Assign & 1.143e4  \\ \hline
Total & 1.227e4  \\ \hline
\end{tabular}
\caption{Total Execution times for 100 iterations of the key steps of the move kernel at three different optimizations.}
\label{tab:GPUPIC_basetime} 
\end{table}
\end{center}


\noindent \begin{table}
\begin{tabular}{| p{4.0cm} | p{3.5cm} | p{3.5cm} |}
\hline
Component & Atomic-Updates (ms) & Sorted+Reduction (ms) \\ \hline
Particle data read, move, and write & 375 & 468 \\ \hline
Potential Grid Read & 467 & 285 \\ \hline
Charge Assign & 1.143e4 & 542 \\ \hline
Particle List Sort & 0 & 2.305e3 \\ \hline
Total & 1.227e4 & 3600 \\ \hline
\end{tabular}
\caption{Total Execution times for 100 iterations of the key steps of the move kernel at two different optimizations.}
\label{tab:GPUPIC_comparison}
\end{table}

\noindent \begin{table}[h]
\begin{tabular}{| p{4.0cm} | p{3.5cm} | p{2.5cm} | p{4.0cm} |}
\hline
Component & SoA (ms) & AoS (ms) & Speedup (SoA vs AoS) \\ \hline
Particle data read, move, and write & 758 & 955 & 1.26x \\ \hline
Count Particles & 32.7 & 109 & 3.35x \\ \hline
Data Reorder & 346 & 480 & 1.38x \\ \hline
Total CPU run time & 2491 & 3284 & 1.31x \\ \hline
\end{tabular}
\caption{Execution times of main steps for Array of Structures and Structure of Arrays. Count Particles and Data Reorder are steps used for a sorted particle list. Count Particles counts the number of particles in each sub-domain. Data Reorder reorders the particle list data after the binindex / particle ID pair have been sorted by the radix sort.}
\label{tab:struct_compare} 
\end{table}


\noindent \begin{table}[h]
\begin{tabular}{| p{4.0cm} | p{3.5cm} | p{2.5cm} | p{4.0cm} |}
\hline
Component & CPU (ns) & GPU (ns) & Speedup \\ \hline
Sort & 0 & 1.240 & - \\ \hline
Charge Assign & 192.055 & 0.621 & 308x \\ \hline
Charge Assign \& Sort & 192.055 & 1.846 & 104x \\ \hline
Poisson Solve & 38.398 & 1.896 & 21x \\ \hline
Particle Advance & 185.245 & 1.520 & 122x \\ \hline
Total\footnote[1] & 417.519 & 5.919 & 71x \\ \hline
\end{tabular}
\caption{CPU and GPU Runtime comparison for 2 GTX 470's vs an Intel(R) Core i7 930 Test was performed using 2 MPI threads handling 21 million particles each on a $64^3$ grid.}
\label{tab:speedup} 
\end{table}

\noindent \begin{table}[h]
\begin{tabular}{| p{4.0cm} | p{3.5cm} | p{2.5cm} | p{4.0cm} |}
\hline
Component & CPU (ns) & GPU (ns) & Speedup \\ \hline
Sort & 0 & 1.272 & - \\ \hline
Charge Assign & 312.210 & 0.802 & 389x \\ \hline
Charge Assign \& Sort & 312.210 & 2.075 & 150x \\ \hline
Poisson Solve & 637.349 & 5.393 & 118x \\ \hline
Particle Advance & 391.335 & 2.325 & 168x \\ \hline
Total\footnote[1] & 1352.461 & 12.958 & 104x \\ \hline
\end{tabular}
\caption{CPU and GPU Runtime comparison for a GTX 590 vs an Intel(R) Xeon(R) CPU E5420. Test was performed using 2 MPI threads handling 17 million particles each on a $64^3$ grid.}
\label{tab:speedup2} 
\end{table} 




\clearpage
\newpage

%\chapter{System Bandwidth Comparison}


In figures \ref{fig:speedup} and \ref{fig:speedup2} we presented an overview for performance of the GPU version of SCEPTIC3D on two distinct systems. System 1 consisted of an Intel(R) Core i7 930 quad core processor clocked at 2.8GHz and two EVGA GTX 470 graphics cards. System 2 consisted of two Intel(R) Xeon(R) E5420 quad core processors clocked at 2.53GHz and one EVGA GTX 590. In figure \ref{fig:speedup2} the speedup factors for the GPU version vs the CPU version were much higher on system 2. There are two reasons for this. The first reason is that the GPU in system 2, the GTX 590, is more powerful than the two GTX 470's in system 1. The second reason has to do with the difference in memory bandwidth between the two systems.

\noindent \begin{figure}
\begin{center}
\includegraphics[width=6in]{appb/bandwidth_test.pdf}
\end{center}
\caption[System Memory Bandwidth Comparison]{System memory bandwidth comparison for the two test systems used to generate the results in \ref{fig:speedup} and \ref{fig:speedup2}.}
\label{fig:memory_bandwidth_compare} 
\end{figure} 

In figure \ref{fig:memory_bandwidth_compare} we performed four vector operations on 20 million element arrays. These vector operations consisted of a simple copy $B = A$, a scalar multiply $B = \alpha A$, a vector add $C = A+B$, and a triad $C = B+\alpha A$. MPI was used to run one copy of the each operation for every core on the target system, the MPI threads were synchronized between operations. Run times were recorded for 50 iterations of the vector operations, the run time at each iteration is the mean of the individual thread run times. The memory bandwidth of each operation is calculated as the number of bytes read and written divided by the minimum run time.    



\clearpage
\newpage

\include{biblio}
\end{document}

